\section{Evaluation}

% In this section, we try to answer these questions:

% 1. How does model's reasoning performance/inference latency/power consumption scale with generation budget (maximum batch size)?

% 2. How does each component contribute to inference latency/resource consumption?

% \sout{3. Can our proposed system generalize to other tasks, e.g., speculative decoding?}

\subsection{Experiment Setup}

\textbf{Devices.} The experiments on NPU performance are conducted on three Android devices: OnePlus Ace3, OnePlus 12, OnePlus Ace5 Pro. Some of the accuracy results are obtained on a server testbed equipped with NVIDIA RTX3090 GPUs.

\begin{table}[]
    \centering
    \begin{tabular}{ccc}
        \hline
        Device & SoC & NPU Arch.\\
        \hline
        OnePlus Ace3 & Snapdragon 8 Gen 2 & V73 \\
        OnePlus 12 & Snapdragon 8 Gen 3 & V75 \\
        OnePlus Ace5 Pro & Snapdragon 8 Elite & V79 \\
        \hline
    \end{tabular}
    \caption{Mobile devices used in evaluation.}
    \label{tab:mobile-devices}
    \vspace{-1em}
\end{table}

\textbf{Models.} We choose models from the Qwen 2.5~\cite{qwen25} and Llama 3.2~\cite{llama32} model family. Considering the actual resource limitations of mobile phones, we mainly evaluate Qwen 2.5 with model sizes of 1.5B and 3B, as well as Llama 3.2 with model sizes of 1B and 3B, which correspond to practical deployable model sizes. When evaluating the performance-cost trade-off of time-time scaling methods, we additionally consider Qwen 2.5 with a model size of 7B. In the evaluation of mathematical reasoning tasks, we use the Instruct model variants of Qwen 2.5 and Llama 3.2. For Best-of-N search and step-level beam search, Skywork-1.5B-PRM~\cite{Skywork-o1} is used as the outcome-reward and process-reward scorer.

\textbf{Datasets and metrics.} In the test-time scaling tasks, we evaluate the pass@1 accuracy of the models in two mathematical reasoning tasks, MATH500~\cite{math} and GSM8K~\cite{gsm8k}, and we uniformly use the 0-shot CoT prompt. For other accuracy measurements, the WinoGrande~\cite{winogrande} accuracy, the MMLU~\cite{mmlu} accuracy, and the Wikitext-2 perplexity are evaluated using \texttt{llama-perplexity} utility.

\textbf{Baselines.} Since we focus on test-time scaling tasks, we mainly present the performance of our implementation under different decoding workloads. To demonstrate the advantages of using NPUs to run test-time scaling workloads, we select the recent OpenCL backend of llama.cpp\footnote{commit: 1caae7f} as the GPU-based system for comparison. This OpenCL backend incorporates optimized Q4\_0 matrix multiplication kernels tailored for Snapdragon's Adreno GPU. Since existing NPU-based systems all have certain limitations in handling test-time scaling workloads, we do not use them as the primary baselines: \texttt{llm.npu}~\cite{mllm-npu} does not utilize the NPU for computation during the decoding phase; other QNN-based systems have low accuracy (e.g., PowerServe~\cite{powerserve}); and systems like Powerinfer-2~\cite{powerinfer2} and HeteroLLM~\cite{heterollm} are not open-source. Nevertheless, we still report the QNN-based data as a reference in Section~\ref{sec:system-comparison}.

% Nevertheless, we have not noticed any work that is quite similar to our setup. \texttt{llm.npu}~\cite{mllm-npu} only uses the Hexagon NPU during the prefilling stage, and the decoding stage is still performed on the CPU. PowerServe~\cite{powerserve} is an open-source project based on QNN. Although its measured speed is consistent with the speed claimed by Qualcomm officially, our measured results show that the model it uses has a significant performance loss in mathematical reasoning tasks (as described previously) and thus fails to meet the most basic accuracy requirements. PowerInfer-2~\cite{powerinfer2} does not have publicly available source code.

\textbf{Settings.} In the operator-level evaluation of GEMM, we select the sizes of the weight matrices of the linear layers corresponding to Qwen2.5-1.5B, Qwen2.5-3B, Llama3.2-1B, and Llama3.2-3B. Specifically, these include the Attention projection matrices $W_q,W_o$ and $W_{gate}, W_{up}, W_{down}$ in the Feed Forward Network (FFN). (For modern models that use Grouped Query Attention (GQA), the projection matrices $W_k,W_v$ in Attention are not selected because their scale is smaller compared to $W_q,W_o$). Most of the matrices adopt the Q4\_0 quantization scheme, which corresponds to 4.5 Bits Per Weight (BPW). As for the FFN down matrices, we apply the Q8\_0 quantization scheme (8.5 BPW) to reduce quantization errors, as existing work indicates their importance in preserving model accuracy~\cite{kovaleva2024systematic,duquant,lin2023emergent}.

\subsection{Overall Performance}

\subsubsection{Accuracy-Latency Trade-off of Test-time Scaling}

\begin{figure*}
    \centering
    \includegraphics[width=1.0\linewidth]{figures/accuracy-latency-tradeoff-v2.pdf}
    \caption{Accuracy-latency trade-off of different test-time scaling methods on various combinations of dataset and hardware.}
    \label{fig:accuracy-latency-tradeoff}
\end{figure*}

Figure~\ref{fig:accuracy-latency-tradeoff} illustrates the performance-cost trade-off of the test-time scaling methods. We use the accuracy in MATH500 and GSM8K as metrics for generation quality and the average decoding latency of on-device models as the cost metric (the data here account for the increased context length introduced by TTS). In the figure, the top row and the bottom row correspond to Best-of-N and Beam Search results, respectively, while "QN"/"LN" denotes the Qwen2.5 or Llama3.2 models with $N$ billion parameters. The SoC results exclude the "8G2" entry due to a known NPU virtual address space limitation~\cite{fastrpc-issue-137} of Snapdragon 8 Gen 2 that prevents models with 3B or more parameters from running. The isolated points marked with a "base" represent the average performance obtained via conventional sampling with the models.

The data show that test-time scaling offers a trade-off space and achieves a more superior Pareto frontier under specific configurations, enabling a better performance-cost balance. In the Best-of-N method, the scaling results of Qwen2.5 1.5B and 3B outperform the baseline accuracies of the 3B and 7B models, respectively. For Beam Search, Qwen2.5-1.5B and Llama3.2-1B can achieve efficiency comparable to or slightly better than their respective 3B variants. Our results indicate that by leveraging the computing power of NPUs and test-time scaling algorithms, small on-device models have the potential to surpass larger models in terms of both generation quality and inference cost.

\subsubsection{On-Device Decoding Performance}

\begin{figure*}[h]
    \centering
    \includegraphics[width=0.7\linewidth]{figures/decoding-throughput.pdf}
    \caption{End-to-End decoding throughput of different models under various batch sizes and hardware settings.}
    \label{fig:decoding-throughput}
\end{figure*}

Figure~\ref{fig:decoding-throughput} demonstrates the on-device decoding throughput of different models in different batch sizes. We only evaluate Qwen2.5-1.5B and Llama-3.2-1B on OnePlus Ace3 due to a 2GiB limitation of the virtual address space on older NPUs.

The data show that for the three devices, the end-to-end decoding throughput of the system significantly increases as the batch size increases. The fundamental reason for the increase in decoding throughput is that the idle computing power of the HMX unit is utilized, and, essentially, the computation time consumed on the core HMX does not increase at all. However, the decoding throughput does not scale perfectly linearly because the inference process contains parts that become much slower with the growth of the input length. Specifically, in our implementation, we conservatively place the weights of the \texttt{lm\_head} (the projection matrix from the hidden states to the vocabulary) and the related activations on the CPU instead of the NPU. Modern LLMs have a large vocabulary, making the \texttt{lm\_head} and logits occupy a large space. Unfortunately, the Hexagon NPU only has a 32-bit virtual address space, therefore placing the complete logits tensor on the NPU may prevent the complete model from running. Currently, we observe that when the batch size equals 16, the proportion of the computation time of logits on the CPU is close to or exceeds 50\%. We expect that after addressing the limitations of the NPU address space and placing the logits computation on the NPU, the system will achieve better throughput scaling characteristics.

\subsubsection{Power and Energy Consumption}

\begin{figure}[h]
    \centering
    \includegraphics[width=1.0\linewidth]{figures/power-consumption.pdf}
    \caption{Power and energy consumption during the LLM decoding stage.}
    \label{fig:power-consumption}
\end{figure}

We measure the power consumption during LLM decoding via sysfs interface on OnePlus 12 with the performance mode enabled. As the batch size increases in the decoding phase, the power consumption of running the 1.5B Qwen model increases, but the overall power consumption of the device is still within 5W; in contrast, the power consumption corresponding to running the 3B Qwen model stabilizes at around 4.3W. Figure~\ref{fig:power-consumption} shows the normalized energy consumption, which is calculated by multiplying the corresponding power consumption by the relative decoding latency. The scaling trait of energy consumption with respect to the batch size is similar to that of decoding latency; therefore, replacing the cost metric in Figure~\ref{fig:accuracy-latency-tradeoff} with energy also results in similar accuracy-cost trade-off characteristics. In particular, we note that the decoding energy consumption of the 1.5B model at a batch size of 8 is lower than that of the 3B model at a batch size of 1, while the test-time scaling accuracy of the 1.5B model when decoding with a batch size of 8 on mathematical tasks is comparable to the base accuracy of the 3B model.

\subsubsection{Comparison with Other Systems} \label{sec:system-comparison}

\begin{figure}[h]
    \centering
    \includegraphics[width=1.0\linewidth]{figures/system-comparison.pdf}
    \caption{Inference throughput comparison.}
    \label{fig:system-comparison}
\end{figure}

The decoding and prefilling performance of our system is presented in Figure~\ref{fig:system-comparison}. We compare our system against a GPU-based implementation and add the performance of FP16 QNN as a reference. During the decoding phase, although the GPU decodes faster at batch size 1, our NPU-based system exhibits higher decoding throughput and better scaling characteristics at larger batch sizes, highlighting the advantage of using NPUs in test-time scaling workloads. Our system also consistently outperforms the GPU-based system in terms of prefilling throughput, achieving comparable performance with proprietary QNN under certain workloads.

\subsection{Accuracy Assessment}

\begin{table}[h]
    \centering
    \begin{tabular}{cccc}
        \hline
        dataset & Tile group & Common group & F16 \\
        \hline
        WinoGrande ($\uparrow$) & 62.559 & 63.349 & 64.613 \\
        MMLU ($\uparrow$) & 35.465 & 35.271 & 34.819 \\
        Wiki PPL ($\downarrow$) & 10.206 & 10.190 & 9.798 \\
        \hline
    \end{tabular}
    \caption{Accuracy comparison between models using tile quantization groups tailored for HMX layout and models using conventional quantization groups.}
    \label{tab:quantization-group-accuracy-compare}
    \vspace{-1em}
\end{table}

\paragraph{Quantization Scheme.} We evaluate the accuracies of the Qwen2.5-1.5B model corresponding to the tile quantization groups based on the HMX layout and the conventional quantization groups. As shown in Table~\ref{tab:quantization-group-accuracy-compare}, the model using our quantization layout has slightly higher accuracy in MMLU compared to the model with the conventional layout, and there is only a slight decrease in Winogrande and Wikitext PPL. Moreover, these accuracy differences are much smaller than the performance loss caused by quantization itself (as indicated by Wikitext perplexity in the "F16" column). In general, using our proposed tile quantization group does not lead to a significant decrease in the accuracy of the quantized model.

\begin{table}[h]
    \centering
    \begin{tabular}{cccc}
        \hline
        dataset & Our LUT16 FA & F32 Attention  \\
        \hline
        WinoGrande ($\uparrow$) & 62.796 & 62.559 \\
        MMLU ($\uparrow$) & 35.207 & 35.465 \\
        Wiki PPL ($\downarrow$) & 10.205 & 10.206 \\
        \hline
    \end{tabular}
    \caption{Accuracy comparison between models using our F16 FlashAttention with LUT-based Softmax and models using conventional F32 Attention.}
    \label{tab:fa-method-accuracy-compare}
    \vspace{-1em}
\end{table}

\paragraph{Attention Implementation.} Table~\ref{tab:fa-method-accuracy-compare} shows the model accuracies corresponding to our LUT-based FP16 Attention and the conventional FP32 Attention, using the same model and datasets as above. It can be seen that replacing the non-critical parts in Attention (except for the accumulation) with a lower FP16 precision does not have a noticeable impact on the end-to-end accuracy of the model.

\subsection{Ablation Study}

\begin{figure}[htbp]
    \centering
    \includegraphics[width=1.0\linewidth]{figures/softmax-exp-compare.pdf}
    \caption{Ablation study of on-chip softmax of our proposed F16 Attention with LUT-based exponential computation. Performance is measured on OnePlus 12.}
    \label{fig:softmax-exp-compare}
\end{figure}

\paragraph{Softmax in Attention.} Figure~\ref{fig:softmax-exp-compare} shows the on-chip softmax latency corresponding to the calculation of the exponential function $\exp$ using different methods under different attention workloads. The length of the input query for Attention is set to 1, 4, and 16, while the length of KV is set to 1024, 4096, or 16384. The figure indicates that our LUT-based exponential calculation achieves an acceleration of 1.26 to 2.19 times compared to the conventional 32-bit floating-point $\exp$, and up to $1.60\times$ speedup compared to the 16-bit floating-point $\exp$. It is worth noting that when pre-computing the $\exp$ lookup table, floating-point numbers with a width of 32 bits or higher can be used to calculate the intermediate results. Therefore, the LUT-based $\exp$ has a higher accuracy than the 16-bit polynomial approximation of $\exp$. When the context length is short, a larger input query will slightly reduce the acceleration ratio, but this phenomenon will be alleviated when the KV length is longer.

\begin{figure*}[htbp]
    \centering
    \includegraphics[width=1.0\linewidth]{figures/gemv-latency-ablation.pdf}
    \caption{Ablation study of proposed optimizations on GEMM dequantization. We measure the performance of GEMV on OnePlus 12.}
    \label{fig:gemv-latency-ablation}
\end{figure*}

\paragraph{Dequantization-based GEMM} Figure~\ref{fig:gemv-latency-ablation} presents the ablation experiment for optimization of the GEMM dequantization layout. The baseline method corresponds to the conventional memory layout, where the column-major weight matrix is quantized according to the continuous groups in memory. The GEMM kernel dequantizes the 32-sized groups one by one during runtime and then scatters the elements to the correct positions in the TCM. The item of "HMX layout" applies the offline weight rearrangement and tile quantization group for the HMX layout, enabling the FP16 weights to be continuously written into the TCM. "Ours" is the version that adopts all the optimizations including the quantization group coalesce. In addition, we add a set of data labeled "no dequantization". In this implementation, instead of performing actual weight dequantization, the quantized weights are read directly from the memory and copied to the on-chip memory without any computation. This set of data can be regarded as the performance upper bound of dequantization-based methods.

Compared to the baseline, our method achieves an acceleration of 9.65 to 19.04 times under different matrix sizes. This is mainly because the scatter operations in the baseline are extremely costly. After applying the HMX layout, the quantization group coalesces and the rearrangements also effectively reduce computational waste, bringing a speedup of $1.82\times$ to $3.45\times$. In particular, compared to the "no dequantization" group, our method is only 27\% slower on average, indicating that this implementation is already close to the performance upper limit of dequantization.

\subsection{Overhead and Sensitivity Analysis}

\begin{figure}[h]
    \centering
    \includegraphics[width=1.0\linewidth]{figures/cpu-memory-usage.pdf}
    \caption{CPU and memory usage during the decoding stage.}
    \label{fig:cpu-memory-usage}
\end{figure}

\paragraph{CPU and Memory Usage.} We evaluate the CPU utilization and memory consumption of the 1.5B and 3B Qwen2.5 models during the decoding stage on OnePlus 12. The CPU memory usage presented in Figure~\ref{fig:cpu-memory-usage} is derived from the resident memory size reported by the \texttt{top} command. We also measure the total size of \texttt{dmabuf} (i.e., memory used by NPU) using \texttt{pmap}, yielding constant values of 1056 MiB and 2090 MiB under a context budget of 4096 tokens for the 1.5B and 3B models, respectively. The total memory consumption is approximately 1.3 GiB for the 1.5B model and 2.4 GiB for the 3B model. The CPU utilization increases with batch size due to the increased computation of vocabulary projection on CPU, yet the number of utilized cores is consistently limited to 4.

\begin{figure}[h]
    \centering
    \includegraphics[width=1.0\linewidth]{figures/prompt-length-sensitivity.pdf}
    \caption{Impact of prompt length on decoding throughput.}
    \label{fig:prompt-length-sensitivity}
\end{figure}

\paragraph{Impact of Prompt Lengths.} Figure~\ref{fig:prompt-length-sensitivity} shows the impact of prompt lengths on decoding throughput. Across all batch sizes and both models, the decoding throughput exhibits a mild decreasing trend as the prompt length increases from 512 to 4096 tokens. However, within the range of prompt lengths up to 4096 tokens, this decline remains relatively subtle, indicating that prompt length exerts only a limited influence on decoding throughput in this interval.
